% --- Archon physics formalization source begin ---
% archon:physics
% archon:covers PhyXMiniProblems/problem_phyx_mini_0872.lean
% archon:source-report reports/phyx_mini/problem_phyx_mini_0872.source.json

\chapter{Physics problem phyx\_mini\_0872}
\label{ch:PhyXMiniProblems_problem_phyx_mini_0872}

\paragraph{Problem source.}
\#\# Question

What is the potential difference \textbackslash{}( \textbackslash{}Delta V\_\{34\} \textbackslash{})?

\#\# Answer choices

- A: A: 12V

- B: B: -14V

- C: C: -200V

- D: D: -20V

\#\# Figure caption (auxiliary; use the image as primary evidence)

The figure is a directed graph with four nodes labeled 1, 2, 3, and 4, connected by arrows. Each arrow represents a voltage change (\textbackslash{}(\textbackslash{}Delta V\textbackslash{})) between nodes:

1. An arrow from node 1 to node 2 with the label \textbackslash{}(\textbackslash{}Delta V\_\{12\} = 30 \textbackslash{}, \textbackslash{}text\{V\}\textbackslash{}).
2. An arrow from node 2 to node 3 with the label \textbackslash{}(\textbackslash{}Delta V\_\{23\} = 50 \textbackslash{}, \textbackslash{}text\{V\}\textbackslash{}).
3. An arrow from node 3 to node 4 with the label \textbackslash{}(\textbackslash{}Delta V\_\{34\}\textbackslash{}).
4. An arrow from node 4 to node 1 with the label \textbackslash{}(\textbackslash{}Delta V\_\{41\} = -60 \textbackslash{}, \textbackslash{}text\{V\}\textbackslash{}).

The arrows represent the direction of voltage change, and the values indicate the magnitude of these changes.

\#\# Dataset metadata

Electric Circuits; Multi-Formula Reasoning, Implicit Condition Reasoning

\paragraph{Recorded answer/context.}
D: D: -20V

\paragraph{Figure/image path.}
/root/proposal\_for\_physic/science-mango/phyx\_mini\_run/phyx\_data/test\_image/872.png

\paragraph{Formalization target.}
create a compiling Lean file with sorry bodies at `PhyXMiniProblems/problem\_phyx\_mini\_0872.lean`. The Lean declarations must preserve the physical quantities, dimensions or dimensional roles, figure labels, governing-law hypotheses, and final relation expressed by this problem.
Use Mathlib/Physlib names found through LeanExplore where available. If a physics API is missing, introduce faithful local abstractions rather than scalar placeholder aliases.

\begin{theorem}[Physics formalization target]
\label{thm:physics:phyx_mini_0872:target}
\lean{PhyXMiniProblems.ProblemPhyXMini0872.deltaV34_eq_negative_twenty_volts_and_matches_choice_D}
\uses{def:physics:phyx-mini-0872:phyxminiproblems-problemphyxmini0872-potentialdifferenceinvolts, def:physics:phyx-mini-0872:phyxminiproblems-problemphyxmini0872-directedvoltageloopsetup, def:physics:phyx-mini-0872:phyxminiproblems-problemphyxmini0872-matchessupplieddirectedvoltagefigure, def:physics:phyx-mini-0872:phyxminiproblems-problemphyxmini0872-satisfiesclosedloopvoltagelaw, def:physics:phyx-mini-0872:phyxminiproblems-problemphyxmini0872-answerchoicepotentialdifferenceinvolts}
The assigned autoformalize agent should translate this physics problem into Lean declarations in the covered file, with theorem and lemma proof bodies written as `by sorry`.
\end{theorem}
\begin{proof}
This is an autoformalization task, not a proof task. Produce statements that can later be proved by the physics prover without weakening the physical model.
\end{proof}

% --- Archon declaration topology begin ---
\section{Declaration topology}
\begin{definition}[energy Dimension]
  \label{def:physics:phyx-mini-0872:phyxminiproblems-problemphyxmini0872-energydimension}
  \lean{PhyXMiniProblems.ProblemPhyXMini0872.energyDimension}
  The physical dimension of energy, M L² T⁻²
\end{definition}
\begin{proof}
  This is the declared model definition of energy Dimension.
\end{proof}
\begin{definition}[electric Potential Difference Dimension]
  \label{def:physics:phyx-mini-0872:phyxminiproblems-problemphyxmini0872-electricpotentialdifferencedimension}
  \lean{PhyXMiniProblems.ProblemPhyXMini0872.electricPotentialDifferenceDimension}
  \uses{def:physics:phyx-mini-0872:phyxminiproblems-problemphyxmini0872-energydimension}
  Electric potential difference has physical dimension energy per charge
\end{definition}
\begin{proof}
  This is the declared model definition of electric Potential Difference Dimension.
\end{proof}
\begin{definition}[Electric Potential Difference Quantity]
  \label{def:physics:phyx-mini-0872:phyxminiproblems-problemphyxmini0872-electricpotentialdifferencequantity}
  \lean{PhyXMiniProblems.ProblemPhyXMini0872.ElectricPotentialDifferenceQuantity}
  \uses{def:physics:phyx-mini-0872:phyxminiproblems-problemphyxmini0872-electricpotentialdifferencedimension}
  A signed, unit-independent physical electric-potential difference
\end{definition}
\begin{proof}
  This is the declared model definition of Electric Potential Difference Quantity.
\end{proof}
\begin{definition}[potential Difference In Volts]
  \label{def:physics:phyx-mini-0872:phyxminiproblems-problemphyxmini0872-potentialdifferenceinvolts}
  \lean{PhyXMiniProblems.ProblemPhyXMini0872.potentialDifferenceInVolts}
  \uses{def:physics:phyx-mini-0872:phyxminiproblems-problemphyxmini0872-electricpotentialdifferencequantity}
  Read a signed electric-potential difference in coherent-SI volts
\end{definition}
\begin{proof}
  This is the declared model definition of potential Difference In Volts.
\end{proof}
\begin{definition}[Circuit Node]
  \label{def:physics:phyx-mini-0872:phyxminiproblems-problemphyxmini0872-circuitnode}
  \lean{PhyXMiniProblems.ProblemPhyXMini0872.CircuitNode}
  The four numbered nodes in image 872.png
\end{definition}
\begin{proof}
  This is the declared model definition of Circuit Node.
\end{proof}
\begin{definition}[Directed Voltage Edge]
  \label{def:physics:phyx-mini-0872:phyxminiproblems-problemphyxmini0872-directedvoltageedge}
  \lean{PhyXMiniProblems.ProblemPhyXMini0872.DirectedVoltageEdge}
  The four potential-change arrows in their displayed cyclic order
\end{definition}
\begin{proof}
  This is the declared model definition of Directed Voltage Edge.
\end{proof}
\begin{definition}[initial Node]
  \label{def:physics:phyx-mini-0872:phyxminiproblems-problemphyxmini0872-directedvoltageedge-initialnode}
  \lean{PhyXMiniProblems.ProblemPhyXMini0872.DirectedVoltageEdge.initialNode}
  \uses{def:physics:phyx-mini-0872:phyxminiproblems-problemphyxmini0872-circuitnode, def:physics:phyx-mini-0872:phyxminiproblems-problemphyxmini0872-directedvoltageedge}
  Initial node of each arrow, as indicated by its voltage subscript
\end{definition}
\begin{proof}
  This is the declared model definition of initial Node.
\end{proof}
\begin{definition}[final Node]
  \label{def:physics:phyx-mini-0872:phyxminiproblems-problemphyxmini0872-directedvoltageedge-finalnode}
  \lean{PhyXMiniProblems.ProblemPhyXMini0872.DirectedVoltageEdge.finalNode}
  \uses{def:physics:phyx-mini-0872:phyxminiproblems-problemphyxmini0872-circuitnode, def:physics:phyx-mini-0872:phyxminiproblems-problemphyxmini0872-directedvoltageedge}
  Final node of each arrow, as indicated by its voltage subscript
\end{definition}
\begin{proof}
  This is the declared model definition of final Node.
\end{proof}
\begin{definition}[printed Label]
  \label{def:physics:phyx-mini-0872:phyxminiproblems-problemphyxmini0872-circuitnode-printedlabel}
  \lean{PhyXMiniProblems.ProblemPhyXMini0872.CircuitNode.printedLabel}
  \uses{def:physics:phyx-mini-0872:phyxminiproblems-problemphyxmini0872-circuitnode}
  Printed numeral identifying each node
\end{definition}
\begin{proof}
  This is the declared model definition of printed Label.
\end{proof}
\begin{definition}[printed Delta VLabel]
  \label{def:physics:phyx-mini-0872:phyxminiproblems-problemphyxmini0872-directedvoltageedge-printeddeltavlabel}
  \lean{PhyXMiniProblems.ProblemPhyXMini0872.DirectedVoltageEdge.printedDeltaVLabel}
  \uses{def:physics:phyx-mini-0872:phyxminiproblems-problemphyxmini0872-directedvoltageedge}
  Printed symbolic label attached to each directed potential change
\end{definition}
\begin{proof}
  This is the declared model definition of printed Delta VLabel.
\end{proof}
\begin{definition}[Figure Node Placement]
  \label{def:physics:phyx-mini-0872:phyxminiproblems-problemphyxmini0872-figurenodeplacement}
  \lean{PhyXMiniProblems.ProblemPhyXMini0872.FigureNodePlacement}
  Coarse relative placement of a node in the supplied graph
\end{definition}
\begin{proof}
  This is the declared model definition of Figure Node Placement.
\end{proof}
\begin{definition}[figure Node Placement]
  \label{def:physics:phyx-mini-0872:phyxminiproblems-problemphyxmini0872-figurenodeplacement-value}
  \lean{PhyXMiniProblems.ProblemPhyXMini0872.figureNodePlacement}
  \uses{def:physics:phyx-mini-0872:phyxminiproblems-problemphyxmini0872-circuitnode, def:physics:phyx-mini-0872:phyxminiproblems-problemphyxmini0872-figurenodeplacement}
  Placement visible in the supplied raster
\end{definition}
\begin{proof}
  This is the declared model definition of figure Node Placement.
\end{proof}
\begin{definition}[Directed Voltage Loop Figure]
  \label{def:physics:phyx-mini-0872:phyxminiproblems-problemphyxmini0872-directedvoltageloopfigure}
  \lean{PhyXMiniProblems.ProblemPhyXMini0872.DirectedVoltageLoopFigure}
  \uses{def:physics:phyx-mini-0872:phyxminiproblems-problemphyxmini0872-circuitnode, def:physics:phyx-mini-0872:phyxminiproblems-problemphyxmini0872-directedvoltageedge, def:physics:phyx-mini-0872:phyxminiproblems-problemphyxmini0872-figurenodeplacement}
  Literal diagram data. printedVoltageInVolts edge = none means that the symbolic ΔV label is present but the raster supplies no numerical value
\end{definition}
\begin{proof}
  This is the declared model definition of Directed Voltage Loop Figure.
\end{proof}
\begin{definition}[Directed Voltage Loop Setup]
  \label{def:physics:phyx-mini-0872:phyxminiproblems-problemphyxmini0872-directedvoltageloopsetup}
  \lean{PhyXMiniProblems.ProblemPhyXMini0872.DirectedVoltageLoopSetup}
  \uses{def:physics:phyx-mini-0872:phyxminiproblems-problemphyxmini0872-electricpotentialdifferencequantity, def:physics:phyx-mini-0872:phyxminiproblems-problemphyxmini0872-directedvoltageedge, def:physics:phyx-mini-0872:phyxminiproblems-problemphyxmini0872-directedvoltageloopfigure}
  The four independent physical potential changes and their associated figure. In particular, the 3 → 4 change is not defined from an answer choice
\end{definition}
\begin{proof}
  This is the declared model definition of Directed Voltage Loop Setup.
\end{proof}
\begin{definition}[Matches Supplied Directed Voltage Figure]
  \label{def:physics:phyx-mini-0872:phyxminiproblems-problemphyxmini0872-matchessupplieddirectedvoltagefigure}
  \lean{PhyXMiniProblems.ProblemPhyXMini0872.MatchesSuppliedDirectedVoltageFigure}
  \uses{def:physics:phyx-mini-0872:phyxminiproblems-problemphyxmini0872-potentialdifferenceinvolts, def:physics:phyx-mini-0872:phyxminiproblems-problemphyxmini0872-figurenodeplacement-value, def:physics:phyx-mini-0872:phyxminiproblems-problemphyxmini0872-directedvoltageloopsetup}
  Facts transcribed from 872.png. The final field connects only those physical potential changes for which the raster actually prints a number. Since the 3 → 4 label has value none, this structure does not constrain the requested potential change numerically
\end{definition}
\begin{proof}
  This is the declared model definition of Matches Supplied Directed Voltage Figure.
\end{proof}
\begin{definition}[Satisfies Closed Loop Voltage Law]
  \label{def:physics:phyx-mini-0872:phyxminiproblems-problemphyxmini0872-satisfiesclosedloopvoltagelaw}
  \lean{PhyXMiniProblems.ProblemPhyXMini0872.SatisfiesClosedLoopVoltageLaw}
  \uses{def:physics:phyx-mini-0872:phyxminiproblems-problemphyxmini0872-potentialdifferenceinvolts, def:physics:phyx-mini-0872:phyxminiproblems-problemphyxmini0872-directedvoltageloopsetup}
  Kirchhoff's voltage law in the orientation of the displayed closed loop. It states the general closed-loop sum relation and does not assign a value to the unknown 3 → 4 edge
\end{definition}
\begin{proof}
  This is the declared model definition of Satisfies Closed Loop Voltage Law.
\end{proof}
\begin{lemma}[delta V34 eq negative sum of other directed changes]
  \label{lem:physics:phyx-mini-0872:phyxminiproblems-problemphyxmini0872-deltav34-eq-negative-sum-of-other-directed-changes}
  \lean{PhyXMiniProblems.ProblemPhyXMini0872.deltaV34_eq_negative_sum_of_other_directed_changes}
  \uses{def:physics:phyx-mini-0872:phyxminiproblems-problemphyxmini0872-potentialdifferenceinvolts, def:physics:phyx-mini-0872:phyxminiproblems-problemphyxmini0872-directedvoltageloopsetup, def:physics:phyx-mini-0872:phyxminiproblems-problemphyxmini0872-satisfiesclosedloopvoltagelaw}
  Rearranging the governing closed-loop law isolates the unknown edge. This is a derived conclusion rather than a field of either premise structure
\end{lemma}
\begin{proof}
  The conclusion follows from the governing relations and figure readouts named in the statement of delta V34 eq negative sum of other directed changes.
\end{proof}
\begin{definition}[Answer Choice]
  \label{def:physics:phyx-mini-0872:phyxminiproblems-problemphyxmini0872-answerchoice}
  \lean{PhyXMiniProblems.ProblemPhyXMini0872.AnswerChoice}
  Labels of the four displayed answer choices
\end{definition}
\begin{proof}
  This is the declared model definition of Answer Choice.
\end{proof}
\begin{definition}[answer Choice Potential Difference In Volts]
  \label{def:physics:phyx-mini-0872:phyxminiproblems-problemphyxmini0872-answerchoicepotentialdifferenceinvolts}
  \lean{PhyXMiniProblems.ProblemPhyXMini0872.answerChoicePotentialDifferenceInVolts}
  \uses{def:physics:phyx-mini-0872:phyxminiproblems-problemphyxmini0872-answerchoice}
  Potential-difference value printed beside each answer, in volts
\end{definition}
\begin{proof}
  This is the declared model definition of answer Choice Potential Difference In Volts.
\end{proof}
% --- Archon declaration topology end ---
% --- Archon physics formalization source end ---
